\documentclass[10pt,a4paper]{amsart}
\usepackage[margin=19mm]{geometry}
\usepackage{amsmath,amssymb,mathtools}
\usepackage[scr=rsfs,cal=euler]{mathalfa}
\usepackage{xfrac}
\usepackage[unicode,hidelinks,bookmarks=false]{hyperref}
\newcommand{\ie}{i.e.\ }
\newcommand{\cf}{cf.\ }
\newcommand{\resp}{respectively\ }
\newcommand{\Subsection}{\subsection}
\providecommand{\nobreakdash}{\nobreak\hbox{-}}
\setlength{\emergencystretch}{2em}
\multlinegap0pt
\usepackage{amssymb}
\usepackage[shortlabels]{enumitem}
\usepackage{csquotes}
\usepackage{tikz}
\usepackage{xcolor}
\newtheorem{lemma}{Lemma}
\newcommand{\X}{\mathbf X}
\newcommand{\Z}{\mathbb{Z}}
\newcommand{\eqdef}{\stackrel{\mathrm{def}}{=}}
\newcommand{\fS}{\mathfrak{S}}
\renewcommand{\ge}{\geqslant}
\renewcommand{\le}{\leqslant}
\renewcommand{\setminus}{\smallsetminus}
\title{A threshold phenomenon for embeddings of Euclidean snowflakes and impossibility of dimension reduction}
\author{Assaf Naor, Kevin Ren}
\date{Source: arXiv:2609.01079v1 (CC BY 4.0)}
\begin{document}
\maketitle

\noindent\emph{Source and licence.} A threshold phenomenon for embeddings of Euclidean snowflakes and impossibility of dimension reduction. Version arXiv:2609.01079v1 (CC BY 4.0), distributed by arXiv under the Creative Commons Attribution 4.0 International licence, http://creativecommons.org/licenses/by/4.0/, as recorded in the arXiv licence metadata for this exact version and shown on the abstract page https://arxiv.org/abs/2609.01079v1. Full licence notice: \texttt{licenses/arxiv-2609.01079-CC-BY-4.0.txt}.\par\smallskip

\noindent\textit{Editorial scope.} Counted unit: the article's lemma fact:abstract (source0.tex lines 449--466). The article's complete original proof of it is delivered with it (source0.tex lines 469--503, one \texttt{proof} environment of 35 lines). No mathematics is written, repaired or re-derived: the statement and the proof are the article's own lines, in the article's own notation, and no retained line is substituted or re-flowed.  No further source-local context is needed: every cross-reference of every retained line resolves inside the two delivered spans.  Nothing was omitted from any retained span, so the package carries no omission record.  No retained line contains a citation, so the wrapper carries no bibliography and no bibliography entry.  No retained line contains a figure environment, an image-inclusion command or drawing code, so no image is delivered and the package declares no asset dependency.  Editorial formatting only, in addition: an AMS \texttt{amsart} preamble that restates the article's own declarations and macros used by the retained lines, namely the environments \texttt{lemma} on separate counters, together with the standard packages those lines need.  The retained environments print in this document as Lemma 1 in order of appearance, whereas the article prints the same items with the numbering of its own full document.  The complete native source of the article as distributed in the arXiv e-print for this exact version is bundled unchanged as \texttt{sources/209078/source0.tex}.
\begin{lemma}\label{fact:abstract} Let  $\fS$ be   a set and suppose that  $L:\fS^2\to [0,\infty)$ satisfies the following requirements:
$$
\forall s,t\in \fS,\qquad L(s,t)=L(t,s)\qquad\mathrm{and}\qquad   L(s,t)=0\iff s=t.
$$
Fix  $0<\beta,\theta\le 1$   and $Q>1$ for which:\footnote{The proof of Lemma~\ref{fact:abstract}  requires $2/3-\theta=\Omega(1)$; the specific upper bound on $\theta$ in~\eqref{Q large theta small} is fixed here for concreteness only.} 
\begin{equation}\label{Q large theta small}
Q\ge \left(\frac{8}{\beta}\right)^{\frac{3}{\theta}}\qquad \mathrm{and}\qquad  \theta\le \frac35. 
\end{equation}
Let $(\X,\|\cdot\|_\X)$ be a Banach space such that for every  $i\in \Z$ there exists a mapping $\phi_i:\fS\to \X$ such that the following estimates hold for every distinct $s,t\in \fS$:
\begin{equation}\label{eq:log version of pairwise}
 \|\phi_i(s)-\phi_i(t)\|_\X\le \min\big\{L(s,t), Q^i\big\}\qquad\mathrm{and}\qquad 
\|\phi_{\lfloor \log_Q L(s,t)\rfloor}(s)-\phi_{\lfloor \log_Q L(s,t)\rfloor}(t)\|_\X\ge \beta Q^{\lfloor \log_Q L(s,t)\rfloor}. 
\end{equation}
Then, there exists a mapping $F:\fS\to \X^3$ that satisfies: 
\begin{equation}\label{eq:F bookkeping guarantees}
\forall s,t\in \fS,\qquad \frac{\beta}{Q^\theta} L(s,t)^{\theta}\lesssim  \|F(s)-F(t)\|_{\X^3}\lesssim Q^\theta L(s,t)^{\theta}. 
\end{equation}
\end{lemma}

\begin{proof} Fix $s_0\in \fS$ and define $f_1,f_2,f_3:\fS\to \X$ as follows: 
\begin{equation}\label{eq:def one of 3 coordinates}
\forall r\in \{1,2,3\},\ \forall s\in \fS,\qquad f_r(s)\eqdef \sum_{m\in \Z} \frac{1}{Q^{(1-\theta)(3m+r)}}\big(\phi_{3m+r}(s)-\phi_{3m+r}(s_0)\big).
\end{equation}
Because~\eqref{Q large theta small} implies that $Q> 3$, the first inequality in~\eqref{eq:log version of pairwise} ensures that the  series  in~\eqref{eq:def one of 3 coordinates} converges absolutely   (at a geometric rate). We can therefore define $F:\fS\to \X^3$ by:
\begin{equation}\label{eq:defF}
\forall s\in \fS,\qquad F(s)\eqdef \big(f_1(s),f_2(s),f_3(s)\big).
\end{equation}



For every distinct $s,t\in  \fS$ let $\lambda_{st}\eqdef \lfloor \log_Q L(s,t)\rfloor$, i.e., $\lambda_{st}$ is the unique integer that satisfies:
\begin{equation}\label{eq:L beteen powers of Q}
Q^{\lambda_{st}}\le L(s,t)<Q^{\lambda_{st}+1}.
\end{equation}
With this notation, the second inequality in~\eqref{eq:log version of pairwise} becomes $\|\phi_i(s)-\phi_i(t)\|_2\ge \beta Q^{\lambda_{st}}$. To demonstrate the first inequality in~\eqref{eq:F bookkeping guarantees}, divide with remainder modulo $3$ to get $\mu_{st}\in \Z$ and $\rho_{st}\in \{1,2,3\}$ satisfying: 
\begin{equation}\label{eq:def mod 3}
\lambda_{st}=3\mu_{st}+\rho_{st}.
\end{equation}
The desired bound is seen by considering as follows the contribution of the $\rho_{st}$ coordinate of $F$: 
\begin{align}\label{eq:lower mod 3}
\begin{split}
\|F(s)-F(t)\|_2&\ge \|f_{\rho_{st}}(s)-f_{\rho_{st}}(t)\|_2\\&\!\!\!\!\!\!\!\!\!\!\!\!\!\!\!\!\!\!\!\stackrel{\eqref{eq:def one of 3 coordinates} \wedge \eqref{eq:defF}\wedge \eqref{eq:def mod 3}}{\ge} \frac{\|\phi_{\lambda_{st}}(s)-\phi_{\lambda_{st}}(t)\|_2}{Q^{(1-\theta)\lambda_{st}}} -\sum_{m\in \Z\setminus \{\mu_{st}\}}
\frac{\|\phi_{\lambda_{st}+3(m-\mu_{st})}(s)-\phi_{\lambda_{st}+3(m-\mu_{st})}(t)\|_2}{Q^{(1-\theta)\left(\lambda_{st}+3(m-\mu_{st})\right)}}\\
&\!\!\!\!\!\!\!\!\!\!\!\stackrel{\eqref{eq:log version of pairwise}\wedge\eqref{eq:L beteen powers of Q}}{\ge} \beta Q^{\theta \lambda_{st}}-\sum_{m=-\infty}^{\mu_{st}-1}Q^{\theta\left(\lambda_{st}+3(m-\mu_{st})\right)}
-\sum^{\infty}_{m=\mu_{st}+1}\frac{Q^{\lambda_{st}+1}}{Q^{(1-\theta)\left(\lambda_{st}+3(m-\mu_{st})\right)}}\\&= \bigg(\beta-\frac{1}{Q^{3\theta}-1}-\frac{Q}{Q^{3(1-\theta)}-1}\bigg)Q^{\theta\lambda_{st}}\stackrel{\eqref{Q large theta small} \wedge \eqref{eq:L beteen powers of Q}}{\gtrsim} \frac{\beta}{Q^\theta} L(s,t)^\theta,
\end{split}
\end{align}
 where the last step of~\eqref{eq:lower mod 3} is a straightforward  calculus exercise. The proof of the rest of~\eqref{eq:F bookkeping guarantees} is simpler: 
\begin{align*}
\|F(s)-F(t)\|_2&\stackrel{\eqref{eq:def one of 3 coordinates} \wedge \eqref{eq:defF}}{\le} \sum_{r=1}^3 \sum_{m\in \Z} \frac{\|\phi_{3m+r}(s)-\phi_{3m+r}(t)\|_2}{Q^{(1-\theta)(3m+r)}}\\
&\stackrel{\eqref{eq:log version of pairwise}\wedge\eqref{eq:L beteen powers of Q}}{\le}  \sum_{i=-\infty}^{\lambda_{st}} Q^{\theta i}+\sum_{i=\lambda_{st}+1}^\infty \frac{Q^{\lambda_{st}+1}}{Q^{(1-\theta)i}} 
=\bigg(\frac{Q^\theta}{Q^\theta-1} +\frac{Q}{Q^{1-\theta}-1}\bigg) Q^{\theta \lambda_{st}}\stackrel{\eqref{Q large theta small}\wedge \eqref{eq:L beteen powers of Q}}{\lesssim} Q^\theta  L(s,t)^{\theta}.\tag*{\qedhere}
\end{align*}
\end{proof}

\end{document}
